\chapter{Introducción}

\section{Presentación del tema y su importancia}

En esta sección se presenta el tema de la monografía, su contexto institucional, académico y práctico, además de la importancia que tiene para el área de Ingeniería de Sistemas.

\section{Planteamiento del problema}

Aquí se describe la situación problemática que motiva el trabajo. Se recomienda delimitar el contexto, identificar las causas principales y explicar las consecuencias que justifican la necesidad del estudio.

\section{Formulación del problema}

La formulación del problema debe expresarse como una pregunta clara, concreta y alineada con el alcance de la monografía.

\section{Objetivo general}

El objetivo general debe expresar el propósito central del trabajo y responder directamente a la formulación del problema.

\section{Objetivos específicos}

\begin{enumerate}[label=\arabic*.]
  \item Definir el primer objetivo específico de la monografía.
  \item Definir el segundo objetivo específico de la monografía.
  \item Definir el tercer objetivo específico de la monografía.
  \item Definir el cuarto objetivo específico de la monografía.
  \item Definir el quinto objetivo específico de la monografía.
\end{enumerate}

\section{Justificación}

La justificación debe desarrollarse en un solo apartado e incluir la relevancia académica, técnica, social o institucional del trabajo, según corresponda.

\section{Alcance y limitaciones}

En esta sección se explica hasta dónde llegará el estudio, qué aspectos serán abordados y cuáles quedarán fuera del análisis.
